\chapter{Rabais, plateformes inversées et paliers}\label{ch:hf:rebates-inverted-venues-and-tiers}

Deux plateformes affichent le même cours acheteur. L'une verse 0.30~cent pour chaque 100 actions exécutées en passif, l'autre en prélève 0.10~cent, si bien que la file est longue sur la première,
courte sur la seconde, et un teneur de marché doit décider où se placer. Dans le modèle de file de ce chapitre, l'ordre qui se place
20\,000 actions derrière sur la plateforme à rabais perd 0.05~cent par action et par ordre envoyé, rabais inclus, car neuf exécutions sur dix qu'il obtient
proviennent du niveau balayé~; l'ordre qui se place 2\,000 actions derrière sur la plateforme inversée atteint le seuil de rentabilité après paiement de ses frais. Le rabais achète une place au fond
d'une longue file, et la file est tarifée en conséquence.

\section{Frais intégrés à la fourchette}

Le livre~1, chapitre~9 décrit la tarification maker-taker et inversée, et le livre~10, chapitre~7 le prix net des frais et le tick effectif~: le
prix économique d'un ordre passif est son prix limite augmenté des frais qu'il acquitte (ou diminué du rabais qu'il perçoit), et des plateformes aux frais différents pour un même
prix affiché sont des prix différents. Pour un teneur de marché, les frais font partie de la fourchette qu'il perçoit.

\begin{definition}[Capture du rebate]\label{def:hf:rebates-inverted-venues-and-tiers:capture}
La \emph{capture du rebate}\index{capture du rebate} (rebate capture) est la part des revenus d'un teneur de marché passif qui provient du rabais versé par la bourse pour l'ajout de liquidité, et non de la
fourchette~; une stratégie qui s'appuie sur elle perçoit le rabais sur des exécutions qui, avant rabais, seraient déficitaires.
\end{definition}

\begin{dated}{2026-09}{Le plafond des frais d'accès et le pilote sur les frais}\label{dat:hf:rebates-inverted-venues-and-tiers:cap}
La Rule 610(c) amendée de la SEC, publiée en octobre 2024, abaisse à 0.001~\$ par action (contre 0.003~\$ fixés en 2005) le plafond des frais pour les exécutions contre des cotations protégées dont le prix est d'au
moins 1.00~\$~; une ordonnance du 31 octobre 2025 a reporté la mise en conformité au premier jour ouvré de novembre
2026. Le Transaction Fee Pilot de la SEC, adopté en décembre 2018 pour tester un plafond plus bas et une interdiction des rabais sur 1\,460 actions tirées au sort, a été
annulé par la cour d'appel du D.C. Circuit le 16 juin 2020.
\end{dated}

Battalio, Corwin et Jennings ont mis en évidence une relation négative entre la qualité d'exécution des ordres à cours limité et le niveau du rabais ou des frais de la plateforme~: les
ordres qui reposent là où les rabais sont élevés s'exécutent moins souvent et subissent davantage de sélection adverse lorsqu'ils le sont. Le modèle de file de la section suivante en montre la raison.

\section{Files sur les plateformes maker-taker et inversées}

Un cours acheteur passif au fond d'une file de $Q$ actions est exécuté par le flux vendeur ordinaire, qui atteint la plateforme au rythme $\lambda$, ou par un balayage
qui prend tout le niveau et fait baisser le prix d'un tick~; ou bien le niveau est abandonné avant l'un ou l'autre, lorsque le cours acheteur s'éloigne. Le premier type d'exécution
rapporte la demi-fourchette, diminuée d'un petit mark-out~; le second perd presque un tick. Une longue file est exécutée surtout par des balayages~: lorsque le flux ordinaire
atteint le fond, le prix a en général déjà bougé.

\begin{omfigure}
\begin{tikzpicture}
\begin{axis}[width=11cm,height=5.2cm,xlabel={\small actions en file devant (milliers)},ylabel={\small edge (cents par action et par ordre)},xmin=0,
  xmax=41,ymin=-0.3,ymax=0.65,tick label style={font=\footnotesize},grid=major,grid style={gray!25},table/col sep=comma,
  legend style={font=\scriptsize,at={(0.97,0.97)},anchor=north east}]
\addplot[omThm,thick,mark=*,mark size=1.3pt] table[x=queue,y=mt]{figdata/hft/17-rebates-inverted-venues-and-tiers/queue.csv};
\addplot[omDef,thick,dashed,mark=square*,mark size=1.3pt] table[x=queue,y=inv]{figdata/hft/17-rebates-inverted-venues-and-tiers/queue.csv};
\draw[gray,dotted] (axis cs:0,0) -- (axis cs:41,0);
\draw[black,thick] (axis cs:20,-0.054) circle[radius=3pt];
\draw[black,thick] (axis cs:2,0.005) circle[radius=3pt];
\legend{{maker-taker : rabais 0.30, 500 actions/s},{inversée : frais 0.10, 300 actions/s}}
\end{axis}
\end{tikzpicture}

\omcaption{Espérance de l'edge par ordre passif de 100 actions, en cents par action, en fonction du nombre d'actions en file devant, sur une plateforme maker-taker (rabais 0.30
cent, flux vendeur de 500 actions par seconde) et une plateforme inversée (frais 0.10 cent, 300 actions par seconde)~; le niveau est balayé toutes les 30 secondes et
abandonné toutes les 20 secondes en moyenne~; une exécution issue du flux ordinaire rapporte un demi-cent diminué de 0.2, une exécution dans un balayage perd un demi-cent~; entourées, les
files typiques (20\,000 et 2\,000 actions). 20\,000 ordres par point.
Données~: \texttt{hf\_rebates.by\_queue}.}\label{fig:hf:rebates-inverted-venues-and-tiers:queue}
\end{omfigure}

À files égales, la plateforme à rabais l'emporte partout (la \cref{fig:hf:rebates-inverted-venues-and-tiers:queue})~: elle verse au lieu de prélever et attire
davantage de flux. Mais les files ne sont pas égales, car chaque teneur de marché a lu la même grille tarifaire. Aux files typiques, 20\,000 actions sur la plateforme à rabais
et 2\,000 sur la plateforme inversée, l'ordre sur la plateforme à rabais s'exécute dans 42~\% des cas, 91~\% de ses exécutions sont des balayages, et il perd $-0.05$ cent par action
et par ordre~; l'ordre sur la plateforme inversée s'exécute dans 74~\% des cas, 24~\% de ses exécutions sont des balayages, et il rapporte $+0.005$. Les files égalisent à peu près les plateformes
comme la concurrence doit le faire, et ce qui reste à un teneur de marché, c'est le timing~: rejoindre une file courte sur la plateforme à rabais (0.40 cent pour 2\,000
actions) est le gain recherché, et il revient à celui qui arrive le premier (chapitre~5).

\section{Paliers et effets de seuil}

Les bourses versent des rabais plus élevés aux membres qui ajoutent davantage de volume, selon des paliers qualifiés par le volume quotidien moyen ajouté au cours du mois, rapporté au
volume consolidé (livre~1, chapitre~29, et livre~10, chapitre~7). Le taux du palier s'applique à chaque action ajoutée dans le mois, ce qui fait de
ce seuil un effet de seuil.

\begin{definition}[Effet de seuil de palier]\label{def:hf:rebates-inverted-venues-and-tiers:cliff}
Un \emph{effet de seuil de palier}\index{effet de seuil de palier} (tier cliff) est un seuil de volume dans une grille tarifaire à partir duquel un meilleur taux s'applique rétroactivement à tout le
volume éligible du mois, de sorte que la facture change de façon discontinue lorsque le seuil est franchi.
\end{definition}

\begin{definition}[Frais marginaux]\label{def:hf:rebates-inverted-venues-and-tiers:marginal}
Les \emph{frais marginaux}\index{frais marginaux} (marginal fee) d'une action sont la variation de la facture totale du membre pour le mois causée par l'ajout de cette action~; en deçà d'un
\omterm{def:hf:rebates-inverted-venues-and-tiers:cliff}{effet de seuil de palier}, ils incluent la valeur du report de tout le volume du mois vers le meilleur palier.
\end{definition}

Sur une grille qui verse 0.20~cent en deçà de 0.20~\% du volume consolidé (11 milliards d'actions par jour) et 0.29~cent au-delà, un membre qui ajoute 18 millions d'actions
par jour paie des \omterm{def:hf:rebates-inverted-venues-and-tiers:marginal}{frais marginaux} de $-0.20$ cent par action~; à 21.9 millions, les 100\,000 actions quotidiennes suivantes valent $-20$ cents chacune, car elles font passer
tout le mois au meilleur palier~; au-delà du seuil, les \omterm{def:hf:rebates-inverted-venues-and-tiers:marginal}{frais marginaux} sont de $-0.29$.

\section{Frais marginaux et décision de volume}

Un membre en deçà d'un seuil peut ajouter un volume qu'il n'aurait pas échangé autrement pour atteindre le palier. Ce volume supplémentaire a un coût~: franchir la fourchette, subir une
sélection adverse, le risque de transactions qu'il ne voulait pas. Avec un coût de 0.6 cent par action de remplissage, un membre dont le volume naturel est de 18 millions d'actions par jour a besoin de
4 millions de plus par jour sur le mois~; le palier vaut 584\,000~\$ en rabais et le remplissage coûte 504\,000~\$~: il est rentable, de 80\,000~\$. À 14 millions de volume naturel, il devrait
en ajouter 8 millions de plus par jour et perdrait 256\,000~\$. Le seuil de rentabilité est de 17.05 millions par jour
(le \cref{lst:hf:rebates-inverted-venues-and-tiers:chase}, la \cref{fig:hf:rebates-inverted-venues-and-tiers:chase})~: la valeur du palier baisse avec
le volume naturel plus lentement que le coût du remplissage. En fin de mois, le supplément nécessaire chaque jour augmente~: si le marché ne peut absorber au plus que 0.1~\% du
volume consolidé par jour venant du membre, la poursuite doit commencer dès le quatorzième des vingt et un jours de bourse.

\begin{omfigure}
\begin{tikzpicture}
\begin{axis}[width=11cm,height=5.4cm,xlabel={\small volume naturel ajouté (millions d'actions par jour)},ylabel={\small milliers de dollars sur le mois},
  xmin=11.5,xmax=22.5,ymin=-500,ymax=1350,tick label style={font=\footnotesize},grid=major,grid style={gray!25},table/col sep=comma,
  legend style={font=\scriptsize,at={(0.5,-0.3)},anchor=north},legend columns=3]
\addplot[omThm,thick,mark=*,mark size=1.3pt] table[x=natural,y=gain]{figdata/hft/17-rebates-inverted-venues-and-tiers/chase.csv};
\addplot[gray,thick,dashed,mark=square*,mark size=1.2pt] table[x=natural,y=cost]{figdata/hft/17-rebates-inverted-venues-and-tiers/chase.csv};
\addplot[omDef,very thick,mark=triangle*,mark size=1.6pt] table[x=natural,y=net]{figdata/hft/17-rebates-inverted-venues-and-tiers/chase.csv};
\draw[black,dotted] (axis cs:11.5,0) -- (axis cs:22.5,0);
\draw[black,dashed] (axis cs:17.05,-500) -- (axis cs:17.05,1350) node[pos=0.93,right,font=\scriptsize] {seuil de rentabilité 17.05};
\draw[black,dashed] (axis cs:22,-500) -- (axis cs:22,1350) node[pos=0.82,left,font=\scriptsize,align=right] {palier\\atteint};
\legend{valeur du palier,coût du remplissage,net}
\end{axis}
\end{tikzpicture}

\omcaption{Poursuite d'un palier à 0.20~\% de 11 milliards d'actions de volume consolidé (22 millions par jour)~: la valeur du meilleur taux sur le
volume du mois, le coût du remplissage à 0.6 cent par action de remplissage et le net, en fonction du volume naturel ajouté du membre, à partir du premier jour. La
poursuite rapporte au-delà de 17.05 millions par jour. Données~: \texttt{hf\_rebates.chase\_curve}.}\label{fig:hf:rebates-inverted-venues-and-tiers:chase}
\end{omfigure}

Les paliers font du volume une décision qui a un prix, et ils donnent au membre une raison de trader qui n'a rien à voir avec ses clients ni avec ses modèles. Trader
avec soi-même pour atteindre un palier relève du wash trading (livre~3, chapitre~16)~: c'est interdit, et c'est détecté.

\section{Fiches de stratégie}

\begin{strategyfile}{Placement passif en capture de rebate}\label{strat:hf:rebates-inverted-venues-and-tiers:capture}
\sfield{Qui paie, et pourquoi} La bourse, qui rémunère la liquidité affichée pour attirer les preneurs~; et les preneurs, à travers la fourchette.
\sfield{Instruments et plateformes} Actions liquides sur des plateformes maker-taker~; valeurs à un seul tick où le rebate représente une grande part de la fourchette.
\sfield{Signal} L'edge net des frais de chaque entrée dans une file~: probabilité d'exécution, part des exécutions issues de balayages, le rebate.
\sfield{Dimensionnement et exécution} N'entrer que dans les files assez courtes pour être exécutées par le flux ordinaire~; annuler lorsque la file devant soi s'allonge ou qu'un balayage devient probable.
\sfield{Coûts} Sélection adverse sur les balayages (91~\% des exécutions en fond de file de 20\,000 actions dans le modèle), compensation et frais réglementaires.
\sfield{Comment elle meurt} Files dont le prix intègre le rebate~; le plafond des frais d'accès plus bas (encadré daté) réduisant les rebates avec lui.
\sfield{Horizon, capacité, infrastructure} Quelques secondes~; suivi de la position dans la file (chapitre~5)~; annulation rapide.
\sfield{Backtest honnête} La file devant chaque ordre à son arrivée, y compris les ordres cachés et les ordres au prix milieu~; frais par palier tels qu'ils étaient.
\sfield{Sources} Battalio, Corwin et Jennings (2016)~; le présent chapitre.
\end{strategyfile}

\begin{strategyfile}{Placement en file sur plateforme inversée}\label{strat:hf:rebates-inverted-venues-and-tiers:inverted}
\sfield{Qui paie, et pourquoi} Les preneurs qui routent vers la plateforme qui les rémunère, et atteignent ainsi en premier la file courte de la plateforme inversée.
\sfield{Instruments et plateformes} Plateformes inversées aux côtés des plateformes maker-taker sur les mêmes valeurs.
\sfield{Signal} La probabilité d'exécution plus élevée et la part de balayages plus faible de la file courte, à comparer à ses frais.
\sfield{Dimensionnement et exécution} Se placer là où la file est assez courte~; passer d'une plateforme à l'autre à mesure que les files évoluent~; utiliser la plateforme inversée pour être le premier à un nouveau
niveau de prix.
\sfield{Coûts} Les frais passifs (0.10 cent dans le modèle).
\sfield{Comment elle meurt} Encombrement (crowding)~: une file inversée aussi longue qu'une file maker-taker perd davantage ($-0.24$ cent pour 20\,000 actions).
\sfield{Horizon, capacité, infrastructure} Quelques secondes~; données de file de toutes les plateformes.
\sfield{Backtest honnête} Règles de routage des preneurs, qui décident de la plateforme dont le flux passe en premier.
\sfield{Sources} Livre~10, chapitre~7~; le présent chapitre.
\end{strategyfile}

\begin{strategyfile}{Gestion du volume tenant compte des paliers}\label{strat:hf:rebates-inverted-venues-and-tiers:tier}
\sfield{Qui paie, et pourquoi} La grille de paliers de la bourse~: le meilleur taux sur tout le volume du mois.
\sfield{Instruments et plateformes} Toutes les plateformes à paliers de volume~; le routage du membre entre elles.
\sfield{Signal} Volume depuis le début du mois comparé aux seuils~; les \omterm{def:hf:rebates-inverted-venues-and-tiers:marginal}{frais marginaux} de chaque action supplémentaire.
\sfield{Dimensionnement et exécution} Concentrer le routage sur la plateforme dont le palier suivant est à portée~; n'ajouter du volume que si la valeur du palier dépasse son coût,
et tôt dans le mois.
\sfield{Coûts} Le coût du trading supplémentaire (0.6 cent par action dans l'exemple)~; risque de concentration sur une seule plateforme.
\sfield{Comment elle meurt} Les grilles changent chaque mois~; des plafonds de frais plus bas compriment les paliers.
\sfield{Horizon, capacité, infrastructure} Un mois~; un moteur de frais (\texttt{firm.feesched}).
\sfield{Backtest honnête} Les grilles en vigueur chaque mois, et la part de volume propre au membre.
\sfield{Sources} Livre~1, chapitre~29~; le présent chapitre.
\end{strategyfile}

\section{Tutoriel~: payé pour attendre, facturé pour attendre}\label{tut:hf:rebates-inverted-venues-and-tiers:tut}

\begin{tutorial}
\textbf{Objectif.} Calculer l'edge net des frais d'un ordre passif sur chaque plateforme, et le seuil de rentabilité de la poursuite d'un palier. \textbf{État final~:}
la \cref{fig:hf:rebates-inverted-venues-and-tiers:queue} et les chiffres des deux dernières sections.
\begin{tutsteps}
\item \textbf{Le modèle de file}~: trois horloges exponentielles par ordre.
\omcode{code/firm/rebatemm/firm_rebatemm.py}{36}{52}{Exécuté par le flux, exécuté par un balayage ou abandonné~: l'edge d'un ordre passif, frais de la plateforme inclus.}\label{lst:hf:rebates-inverted-venues-and-tiers:queue}
\item \textbf{Les plateformes} et le balayage selon la longueur des files (\texttt{hf\_rebates.venues}, \texttt{hf\_rebates.by\_queue}).
\item \textbf{Le palier}~: \omterm{def:hf:rebates-inverted-venues-and-tiers:marginal}{frais marginaux} et poursuite sur la grille de \texttt{firm.feesched}.
\omcode{code/firm/rebatemm/firm_rebatemm.py}{64}{77}{La poursuite~: le volume supplémentaire nécessaire, la valeur du palier sur toutes les actions du mois, et son coût.}\label{lst:hf:rebates-inverted-venues-and-tiers:chase}
\end{tutsteps}
\textbf{Pour aller plus loin.} Placer les files sur \texttt{firm.exchsim} avec deux plateformes et un routage réel~; faire dépendre le taux de balayage de la
taille de la file~; ajouter une seconde grille et router entre les deux.
\end{tutorial}

\section{Construction~: le module de rebate et de paliers}\label{bld:hf:rebates-inverted-venues-and-tiers:rebatemm}

\begin{build}
\bfield{Rôle} Pricer un ordre passif par plateforme, frais et file inclus, et décider du volume du mois autour des \omterm{def:hf:rebates-inverted-venues-and-tiers:cliff}{effets de seuil de palier}.
\bfield{Interface} \texttt{queue\_edge(Q, lam, mu\_through, mu\_away, fee, half\_spread, adverse, lot, n, seed)}, \texttt{monthly\_bill},
\texttt{marginal\_fee}, \texttt{chase(schedule, natural, tcv, days\_left, days, loss)}, \texttt{last\_day\_to\_chase}. Construit sur
\texttt{firm.feesched} (livre~1).
\bfield{Règles} Des cents par action~; les frais positifs sont payés par le membre~; les taux de palier s'appliquent à tout le volume ajouté du mois.
\bfield{Tests d'acceptation} \texttt{code/firm/rebatemm/tests/}~: sans balayage, l'edge vaut la demi-fourchette moins le mark-out moins les frais~; les files plus longues
s'exécutent moins et davantage par balayage~; la facture mensuelle et les \omterm{def:hf:rebates-inverted-venues-and-tiers:marginal}{frais marginaux} à la main, le pic du seuil~; les volumes, le gain et le coût de la poursuite à la main~; le module
échoue à faible volume naturel~; le dernier jour pour commencer.
\bfield{Pour aller plus loin} Files sur \texttt{firm.exchsim}~; routage entre plusieurs grilles~; effet de la baisse du plafond des frais d'accès sur la longueur des files.
\end{build}

\begin{omsources}
\item R.~Battalio, S.~A.~Corwin, R.~Jennings, Can brokers have it all? On the relation between make-take fees and limit order execution quality,
\emph{Journal of Finance} 71(5), 2016, 2193--2238.
\item US Securities and Exchange Commission, amendments to Rule 610, \emph{Federal Register} 89 FR 81620, 8 octobre 2024~; communiqué de presse 2025-130.
\item \emph{New York Stock Exchange LLC v. SEC}, No.\ 19-1042 (D.C. Cir.\ 2020).
\end{omsources}

\section{Exercices}

\begin{exercise}[$\star$]\label{exo:hf:rebates-inverted-venues-and-tiers:1}
Un cours acheteur de 20.00~\$ sur une plateforme versant un rebate de 0.30 cent, et un cours acheteur de 20.00~\$ sur une plateforme facturant des frais de 0.10 cent. Calculer leurs
prix nets des frais pour l'acheteur passif.
\end{exercise}

\begin{exercise}[$\star$]\label{exo:hf:rebates-inverted-venues-and-tiers:2}
Un membre ajoute 18 millions d'actions par jour à un rebate de 0.20 cent. Calculer le rebate du mois sur 21 jours, puis à 0.29 cent.
\end{exercise}

\begin{exercise}[$\star$]\label{exo:hf:rebates-inverted-venues-and-tiers:3}
Pourquoi les \omterm{def:hf:rebates-inverted-venues-and-tiers:marginal}{frais marginaux} juste en deçà d'un \omterm{def:hf:rebates-inverted-venues-and-tiers:cliff}{effet de seuil de palier} sont-ils si élevés~?
\end{exercise}

\begin{exercise}[$\star\star$]\label{exo:hf:rebates-inverted-venues-and-tiers:4}
Pourquoi un ordre en fond de longue file ne reçoit-il surtout que des exécutions toxiques~?
\end{exercise}

\begin{exercise}[$\star\star$]\label{exo:hf:rebates-inverted-venues-and-tiers:5}
Établir le volume naturel de seuil de rentabilité de la poursuite d'un palier~: seuil du palier 22 millions, rebate actuel 0.20, rebate suivant 0.29, coût de 0.6 par action de remplissage.
\end{exercise}

\begin{exercise}[$\star\star$]\label{exo:hf:rebates-inverted-venues-and-tiers:6}
Que fait la baisse du plafond des frais d'accès à la file de la plateforme à rebate~?
\end{exercise}

\begin{exercise}[$\star\star\star$]\label{exo:hf:rebates-inverted-venues-and-tiers:7}
\emph{Programmation.} Avec \texttt{firm.rebatemm.queue\_edge}, trouver la longueur de file à laquelle l'edge par ordre de la plateforme à rebate tombe à zéro.
\end{exercise}

\begin{exercise}[$\star\star\star$]\label{exo:hf:rebates-inverted-venues-and-tiers:8}
\emph{Trouver la faille.} «~Le rebate est de 0.30 cent et notre demi-fourchette est de 0.5, donc chaque exécution passive rapporte 0.8 cent.~»
\end{exercise}

\section{Problème~: Payé pour attendre, facturé pour attendre}

\begin{problem}[{Problème du week-end --- payé pour attendre, facturé pour attendre}]\label{pb:hf:rebates-inverted-venues-and-tiers:1}
Un teneur de marché choisit où laisser reposer ses cours acheteurs et quel volume ajouter dans le mois.

\medskip\noindent\textbf{Partie I --- Frais.}
\begin{enumerate}
\item Définir la \omterm{def:hf:rebates-inverted-venues-and-tiers:capture}{capture du rebate} (rebate capture).
\item Résumer l'encadré daté.
\item Que trouvent Battalio, Corwin et Jennings~?
\item Pourquoi un teneur de marché perçoit-il des frais faisant partie de la fourchette~?
\end{enumerate}

\medskip\noindent\textbf{Partie II --- Files.}
\begin{enumerate}[resume]
\item Décrire les trois issues d'un ordre au repos dans le modèle.
\item Donner la probabilité d'exécution, la part de balayage et l'edge aux files typiques.
\item Pourquoi les files égalisent-elles les plateformes~?
\item Où se trouve le gain, et qui le reçoit~?
\end{enumerate}

\medskip\noindent\textbf{Partie III --- Paliers.}
\begin{enumerate}[resume]
\item Définir un \omterm{def:hf:rebates-inverted-venues-and-tiers:cliff}{effet de seuil de palier} (tier cliff) et les \omterm{def:hf:rebates-inverted-venues-and-tiers:marginal}{frais marginaux}.
\item Donner les \omterm{def:hf:rebates-inverted-venues-and-tiers:marginal}{frais marginaux} à 18, 21.9 et 22 millions d'actions par jour.
\item Calculer la poursuite à 18 et 14 millions.
\item Quand la poursuite doit-elle commencer~?
\end{enumerate}

\medskip\noindent\textbf{Partie IV --- Le verdict.}
\begin{enumerate}[resume]
\item Énoncer le \emph{résultat nommé}~: l'edge net des frais de rester au repos sur chaque plateforme, et le volume de fin de mois auquel poursuivre le palier suivant cesse
d'être rentable.
\item Sur quelle plateforme rester au repos, et quand~?
\item Que changerait un plafond des frais d'accès plus bas~?
\item Où se situe la frontière entre poursuivre un palier et le wash trading~?
\item Comment un moteur de frais doit-il alimenter le moteur de cotation~?
\item Quelle fiche de stratégie est la plus exposée aux changements de calendrier~?
\item Comment mesurer la part des exécutions qui sont des balayages~?
\item En une phrase~: que achète un rebate~?
\end{enumerate}
\end{problem}

\section{Questions d'entretien}

\begin{interviewq}[$\star$ \iqroles{trader}]\label{iq:hf:rebates-inverted-venues-and-tiers:1}
Pourquoi quelqu'un paierait-il pour afficher un ordre sur une plateforme inversée~?
\end{interviewq}

\begin{interviewq}[$\star\star$ \iqroles{chercheur}]\label{iq:hf:rebates-inverted-venues-and-tiers:2}
Comment estimeriez-vous la probabilité d'exécution et le mark-out d'un ordre passif en fonction de la file devant lui, à partir de vos propres données d'ordres~?
\end{interviewq}

\begin{interviewq}[$\star\star$ \iqroles{trader}]\label{iq:hf:rebates-inverted-venues-and-tiers:3}
Nous sommes le 19 et vous êtes à 3~\% sous un seuil de palier sur une plateforme. Que faites-vous~?
\end{interviewq}

\begin{interviewq}[$\star\star$ \iqroles{développeur}]\label{iq:hf:rebates-inverted-venues-and-tiers:4}
Concevez un moteur de frais qui fournit au moteur de cotation, en temps réel, les \omterm{def:hf:rebates-inverted-venues-and-tiers:marginal}{frais marginaux} de l'action suivante sur chaque plateforme.
\end{interviewq}

\begin{interviewq}[$\star\star$ \iqroles{risque}]\label{iq:hf:rebates-inverted-venues-and-tiers:5}
Le P\&L d'un desk provient à 80~\% de rebates. Quels risques cela concentre-t-il~?
\end{interviewq}

\begin{interviewq}[$\star\star\star$ \iqroles{chercheur}]\label{iq:hf:rebates-inverted-venues-and-tiers:6}
Dans le modèle à trois horloges, établir la probabilité que l'ordre soit exécuté par le flux avant un balayage ou un abandon lorsque la file est consommée
à vitesse constante.
\end{interviewq}
